\chapter{Đào sâu: JWT ở quy mô lớn --- xác minh ở đâu, token béo phì và khóa bất đối xứng}
\label{ch:jwt-scale}

Chương~5 dạy thiết kế bảo mật của dự án đúng như nó đang có: token
HS256, xác minh một lần tại gateway, danh tính chuyển tiếp dưới dạng
header. Thiết kế đó đúng với một hệ thống có một gateway và năm service.
Chương này là chương \emph{đào sâu} đầu tiên của cuốn sách: lấy một thiết
kế đã có trong dự án và đẩy nó tới khi gãy, theo cách một người phỏng vấn
hay một đợt tăng lưu lượng gấp mười sẽ làm. Câu hỏi lấy từ một bài viết
tiếng Việt được chia sẻ rộng rãi về JWT trong
microservices;\footnote{``Lột trần sự thật về JWT --- Phần 2: Trận chiến
Microservices và nỗi ám ảnh Token béo phì'', viblo.asia, 2026.} câu trả
lời lấy từ việc đọc code của dự án.

Ba câu hỏi sắp xếp chương này:

\begin{enumerate}
  \item \textbf{Nên xác minh token ở đâu} --- một lần tại gateway, hay
  trong từng service?
  \item \textbf{Chuyện gì xảy ra khi token phình to} --- và mẫu
  \emph{phantom token} giữ nó gọn là gì?
  \item \textbf{Chữ ký thực sự bảo đảm điều gì}, và vì sao lựa chọn
  giữa HS256 và RS256 quyết định ai được tin cậy với thứ gì?
\end{enumerate}

\section{Xác minh ở đâu: gateway hay từng service?}
\label{sec:where-to-verify}

\subsection{Mô hình 1: xác minh tập trung (điều dự án đang làm)}

Đọc lại hai nửa của mẫu biên tin cậy. Ở biên, trong
\code{JwtAuthenticationFilter} của gateway:

\begin{lstlisting}[language=Java]
Claims claims = Jwts.parser()
        .verifyWith(signingKey)                    (*@\co{1}@*)
        .build()
        .parseSignedClaims(token)
        .getPayload();

String userId = claims.getSubject();
String role = claims.get(JwtConstants.CLAIM_ROLE, String.class);

ServerHttpRequest mutatedRequest = exchange.getRequest().mutate()
        .header(JwtConstants.HEADER_USER_ID, userId)   (*@\co{2}@*)
        .header(JwtConstants.HEADER_USER_ROLE, role)
        .build();
\end{lstlisting}

Và phía sau nó, trong \code{GatewayHeaderAuthFilter} của mỗi service
(bản của course-service):

\begin{lstlisting}[language=Java]
String userId = request.getHeader("X-User-Id");
String userRole = request.getHeader("X-User-Role");

if (userId != null && userRole != null) {               (*@\co{3}@*)
    var authorities = List.of(
            new SimpleGrantedAuthority("ROLE_" + userRole));
    var authentication = new UsernamePasswordAuthenticationToken(
            Long.valueOf(userId), null, authorities);
    SecurityContextHolder.getContext().setAuthentication(authentication);
}
\end{lstlisting}

\begin{description}
  \item[\co{1}] Mật mã chạy đúng một lần, ở đây. Khóa chính là
  \code{app.jwt.secret} mà auth-service dùng để ký.
  \item[\co{2}] Danh tính rời gateway dưới dạng hai header HTTP thuần.
  Không gì ở phía sau phân biệt được gateway viết chúng hay ai khác
  viết.
  \item[\co{3}] Service \emph{không} kiểm tra gì. Có hai header nghĩa là
  đã xác thực. Đó là toàn bộ hợp đồng.
\end{description}

Ưu điểm là thật: một nơi giữ secret, service không cần thư viện JWT lẫn
xử lý khóa, và chi phí CPU kiểm tra chữ ký mỗi request chỉ trả một lần.
Phần lớn tutorial Spring Cloud dừng ở đây.

\subsection{Ngụy biện ``mạng nội bộ đáng tin''}

Mẫu này dựa trên một bất biến --- \emph{mọi} đường tới một service đều
đi qua gateway. Trong file Compose của dự án, bất biến đó đã sai sẵn.
\code{docker-compose.yml} publish thẳng auth-service ở \code{8081} và
course-service ở \code{8082} ra host, để bạn debug được. Bất kỳ ai trên
host vì thế có thể bỏ qua gateway:

\begin{handson}[giả mạo danh tính trong một dòng]
Với stack đang chạy, hỏi course-service danh sách khóa học của một giáo
viên mà không hề đưa token:
\begin{lstlisting}[language=bash]
$ curl -s -H "X-User-Id: 1" -H "X-User-Role: TEACHER" \
       http://localhost:8082/courses | head -c 300
\end{lstlisting}
Nó trả lời. Giờ lặp lại qua gateway ở \code{8080} với cùng hai header và
không có \code{Authorization}:
\begin{lstlisting}[language=bash]
$ curl -s -H "X-User-Id: 1" -H "X-User-Role: TEACHER" \
       http://localhost:8080/api/courses
{"error":"Unauthorized","message":"Missing or invalid Authorization header"}
\end{lstlisting}
Khác biệt giữa hai kết quả chính là toàn bộ mô hình bảo mật.
\end{handson}

Trong Kubernetes các cổng publish biến mất (Chương~18), nhưng mạng pod
phẳng vẫn còn: pod nào trong namespace cũng với được
\code{course-service:8082}. Ngụy biện nằm ở chỗ coi ``bên trong'' nghĩa
là ``đáng tin''. Chỉ cần một container bị chiếm --- một image có lỗ
hổng, một kubeconfig bị lộ, một dependency trong chuỗi cung ứng --- kẻ
tấn công giả mạo \code{X-User-Id: 1} tới mọi service cùng lúc.

\begin{pitfall}[header danh tính sống sót qua đường dẫn công khai]
Triệu chứng: không thấy gì, và đó chính là vấn đề. Nhìn đầu hàm
\code{JwtAuthenticationFilter.filter}: request tới đường dẫn công khai
(\code{/api/auth/refresh}, \code{/actuator/**}) được chuyển tiếp
\emph{nguyên vẹn} --- gateway không xác minh token cũng không đụng tới
header. Client vì thế có thể gửi \code{X-User-Id: 1} tới một đường dẫn
công khai và header tới service còn nguyên. Hôm nay không endpoint công
khai nào đọc nó, nên chưa khai thác được; nhưng bất biến ``chỉ gateway
mới viết các header này'' không được cưỡng chế, nó chỉ đúng nhờ may
mắn. Cách sửa là một default filter trong \code{api-gateway.yml} gỡ cả
hai header khỏi mọi request \emph{đi vào} trước khi filter JWT đặt
chúng:
\begin{lstlisting}[language=yaml]
spring:
  cloud:
    gateway:
      default-filters:
        - RemoveRequestHeader=X-User-Id
        - RemoveRequestHeader=X-User-Role
\end{lstlisting}
Giờ các header chỉ có thể xuất phát từ gateway. Việc này đóng cánh cửa
bên ngoài; nó không làm gì cho cánh cửa trong cụm.
\end{pitfall}

\subsection{Mô hình 2: zero trust --- service nào cũng xác minh}

Phương án thay thế gỡ bỏ sự tin cậy: mỗi service tự parse và xác minh
JWT, và bỏ qua mọi header danh tính. Header giả giờ vô giá trị vì
service không bao giờ đọc; token giả thì trượt kiểm tra chữ ký. Gateway
vẫn có thể xác minh --- loại token hỏng từ sớm rất rẻ --- nhưng không gì
ở phía sau \emph{phụ thuộc} vào việc đó.

Chi phí xuất hiện ngay: mọi bên xác minh đều cần khóa. Với HS256 nghĩa
là chép \emph{secret ký} tới mọi service, còn tệ hơn vấn đề nó giải
quyết (Mục~\ref{sec:hs-vs-rs}). Zero trust ở tầng ứng dụng chỉ hợp lý
với khóa bất đối xứng, và đó là lý do hai nửa của chương này đi cùng
nhau.

\begin{decision}[header do gateway xác minh, cho tới lúc này]
Dự án giữ Mô hình~1. Nó có đúng hai nơi giữ khóa, không có service bên
thứ ba, và mọi service đều do cùng một pipeline triển khai --- bán kính
thiệt hại của giả định biên tin cậy chỉ là một namespace. Ba tín hiệu
sẽ lật quyết định: một service do đội khác viết, một service mà bên gọi
không chỉ là gateway (một consumer Kafka cũng mở HTTP, một batch job),
hoặc một yêu cầu tuân thủ đòi danh tính phải xác minh được tại điểm sử
dụng. Lúc đó bước đi là RS256 cộng xác minh tại từng service, không
phải thêm header. Hệ thống production thường đi đường thứ ba: giữ
header, và để \emph{tầng vận chuyển} chứng minh người gửi bằng mutual
TLS từ một service mesh (Istio, Linkerd). Gateway khi đó thành một đối
tác được định danh bằng mật mã thay vì một hàng xóm không kiểm chứng
được.
\end{decision}

\section{Chữ ký bảo đảm điều gì}

Một JWT là \code{base64url(header).base64url(payload).signature}. Hai
phần đầu được \emph{encode}, không phải mã hóa: dán bất kỳ token nào của
dự án vào một bộ giải mã và \code{sub}, \code{email}, \code{role} đọc
được ngay. Chữ ký thêm đúng hai tính chất:

\begin{itemize}
  \item \textbf{Toàn vẹn.} Đổi một byte của payload
  (\code{"role":"STUDENT"} thành \code{"TEACHER"}) là xác minh thất bại.
  \item \textbf{Xác thực.} Token do một bên giữ khóa ký tạo ra --- và
  không ai khác.
\end{itemize}

Nó không thêm tính bí mật. Cả HS256 lẫn RS256 đều tính chữ ký trên cùng
một chuỗi, \code{base64url(header) + "." + base64url(payload)}. Chúng
chỉ khác ở \emph{khóa nào tạo chữ ký và khóa nào kiểm tra nó}. Một khác
biệt duy nhất đó quyết định toàn bộ cấu trúc tin cậy của một hệ thống.

\section{HS256 so với RS256}
\label{sec:hs-vs-rs}

\subsection{HS256: một khóa, hai vai}

HS256 là HMAC-SHA256: một hàm băm có trộn khóa. Bạn không cần nhớ cấu
trúc (hai lượt SHA-256 lồng nhau với khóa XOR cùng hai pad cố định); bạn
cần một tính chất. \emph{Cùng một khóa vừa tạo vừa kiểm tra chữ ký.}
Xác minh nghĩa là tính lại HMAC với bản khóa của mình rồi so byte. Không
có khóa thì không tính được, nên không giả mạo được. Có khóa thì tính
được, nên \emph{có thể} giả mạo. Bên xác minh và bên ký không phân biệt
được với nhau.

Hai bản sao của khóa đó trong dự án:

\begin{lstlisting}[language=Java]
// auth-service/.../service/JwtService.java
this.signingKey = Keys.hmacShaKeyFor(secret.getBytes(StandardCharsets.UTF_8));

// api-gateway/.../filter/JwtAuthenticationFilter.java
this.signingKey = Keys.hmacShaKeyFor(jwtSecret.getBytes(StandardCharsets.UTF_8));
\end{lstlisting}

Cả hai đọc \code{app.jwt.secret}; trong cụm cả hai mount cùng một Secret
\code{jwt} (Chương~16). Gateway chỉ bao giờ cần \emph{xác minh}, vậy mà
nó giữ một khóa có thể \emph{đúc} token cho bất kỳ user với bất kỳ vai
trò nào. Chiếm được pod gateway là sở hữu mọi danh tính trong hệ thống.
Mở rộng Mô hình~2 một cách ngây thơ --- chép secret tới năm service ---
và bạn có sáu nơi như vậy.

\subsection{RS256: hai khóa, mỗi khóa một vai}

RS256 là chữ ký RSA trên một mã băm SHA-256. RSA dựa trên một bất đối
xứng trong số học: nhân hai số nguyên tố lớn thì tức thì; tìm lại chúng
từ tích thì bất khả thi.

Sinh khóa, một lần:

\begin{enumerate}
  \item Chọn hai số nguyên tố lớn $p$, $q$ (mỗi số khoảng 1024 bit cho
  khóa 2048 bit). Tính $n = p \times q$.
  \item Chọn số mũ công khai $e$; thực tế luôn là $65537$. Cặp $(n, e)$
  là \textbf{khóa công khai}.
  \item Tính $d$ sao cho $e \times d \equiv 1 \pmod{(p-1)(q-1)}$. Cặp
  $(n, d)$ là \textbf{khóa riêng}. Tính $d$ cần $p$ và $q$; chỉ biết $n$
  thì không suy ra được chúng.
\end{enumerate}

Ký và xác minh:

\begin{lstlisting}
h = SHA256(header "." payload), padded (PKCS#1 v1.5)
sign:    signature = h^d  mod n      -- needs d: only auth-service
verify:  signature^e mod n == h ?    -- needs e, n: anyone
\end{lstlisting}

Vì $e \times d \equiv 1$, nâng lũy thừa $d$ rồi lũy thừa $e$ trả về $h$.
Ai có khóa công khai đều kiểm tra được chữ ký; không ai thiếu $d$ tạo
được chữ ký. Đó là bất đối xứng mà HS256 không có, và nó thay đổi ai
được giữ thứ gì:

\begin{figure}[htb]
\centering
\begin{tikzpicture}[node distance=6mm and 15mm]
  \node[boxdark] (auth) {auth-service\\khóa riêng $(n,d)$};
  \node[boxfill, right=of auth] (jwks) {endpoint JWKS\\khóa công khai $(n,e)$, \code{kid}};
  \node[box, right=of jwks] (gw) {api-gateway\\chỉ xác minh};
  \node[box, above=of gw] (dict) {dictionary-service\\chỉ xác minh};
  \node[box, below=of gw] (course) {course-service\\chỉ xác minh};
  \draw[flow] (auth) -- node[lbl,above]{công bố} (jwks);
  \draw[flowdash] (jwks.east) -- node[lbl,below,pos=0.55]{tải, cache} (gw.west);
  \draw[flowdash] (jwks.east) -- (dict.west);
  \draw[flowdash] (jwks.east) -- (course.west);
\end{tikzpicture}
\caption{Cấu trúc khóa RS256: một bên ký, bao nhiêu bên xác minh cũng
được. Lộ một bên xác minh không lộ gì cả.}
\label{fig:rs256-topology}
\end{figure}

\begin{itemize}
  \item Auth-service giữ khóa riêng và là nơi \emph{duy nhất} có thể
  phát hành token.
  \item Gateway và mọi service giữ khóa công khai. Lộ nó vô hại --- nó
  công khai theo định nghĩa. Bạn có thể in nó trong cuốn sách này.
  \item Phân phối khóa không còn là bài toán quản lý bí mật: công bố
  khóa công khai tại một URL quy ước (một \emph{JWKS}, JSON Web Key
  Set), và bên xác minh nào cũng tải lúc khởi động rồi cache lại.
\end{itemize}

Bạn đã tiêu thụ cấu trúc này rồi. Đăng nhập Google ở Chương~5 trả về
một \code{id\_token} ký RS256; OAuth2 client của Spring xác minh nó với
JWKS của Google tại \code{https://www.googleapis.com/oauth2/v3/certs}.
Google không bao giờ đưa bạn khóa riêng, và bạn không bao giờ cần nó.

\subsection{Chi phí}

\begin{center}
\begin{tabular}{@{}lll@{}}
\toprule
 & HS256 & RS256 (2048-bit) \\
\midrule
Vật liệu khóa & một secret, $\ge$ 32 byte & cặp riêng/công khai, $\approx$1,7\,KB PEM \\
Kích thước chữ ký & 32 byte (43 ký tự) & 256 byte (342 ký tự) \\
Ký & micro giây & khoảng 1\,ms \\
Xác minh & micro giây & vài chục micro giây \\
Ai ký được & bất kỳ ai giữ khóa & chỉ người giữ khóa riêng \\
Ai xác minh được & bất kỳ ai giữ khóa & bất kỳ ai \\
\bottomrule
\end{tabular}
\end{center}

Xác minh RSA rẻ hơn ký rất nhiều vì $e = 65537$ rất nhỏ trong khi $d$
rộng 2048 bit. Bất đối xứng đó khớp với cấu trúc: một bên ký, nhiều bên
xác minh. Nếu kích thước chữ ký quan trọng, ES256 (ECDSA trên đường cong
P-256) cho chữ ký 64 byte và khóa nhỏ hơn nhiều, và là mặc định hiện
đại cho các identity provider mới.

\begin{pitfall}[nhầm lẫn thuật toán (algorithm confusion)]
Triệu chứng: một token đáng lẽ bị từ chối lại được chấp nhận --- và log
không thấy gì sai. Nguyên nhân: thư viện cũ đọc trường \code{alg} trong
header để quyết định \emph{cách} xác minh. Kẻ tấn công lấy một token
RS256 hợp lệ, sửa header thành \code{"alg":"HS256"}, và ký bằng HMAC với
secret là \emph{chính các byte của khóa công khai} --- những byte vốn dĩ
công khai. Bên xác minh ngây thơ thấy HS256, dùng khóa công khai mình có
làm secret HMAC, và phép kiểm tra qua. JJWT chặn việc này bằng cách để
\emph{loại khóa} chọn thuật toán, không phải header:
\code{verifyWith(SecretKey)} từ chối mọi token bất đối xứng và
\code{verifyWith(PublicKey)} từ chối mọi token HMAC. Đó là lý do
\code{JwtAuthenticationFilter} an toàn hôm nay, và là lý do bạn không
bao giờ được xác minh bằng một phương thức chấp nhận ``header nói gì thì
làm nấy''.
\end{pitfall}

\subsection{Chuyển đổi dự án}

Thay đổi đủ nhỏ để đọc trọn vẹn. Sinh cặp khóa một lần:

\begin{lstlisting}[language=bash]
$ openssl genpkey -algorithm RSA -pkeyopt rsa_keygen_bits:2048 \
      -out jwt-private.pem
$ openssl pkey -in jwt-private.pem -pubout -out jwt-public.pem
\end{lstlisting}

PEM riêng chỉ đi vào Secret của auth-service; PEM công khai có thể đặt
trong ConfigMap, vì nó không bí mật. \code{JwtService} ký bằng khóa
riêng:

\begin{lstlisting}[language=Java]
@Service
public class JwtService {

    private final PrivateKey privateKey;   // sign
    private final PublicKey publicKey;     // self-verify where needed

    public JwtService(@Value("${app.jwt.private-key-pem}") String privPem,
                      @Value("${app.jwt.public-key-pem}") String pubPem)
            throws GeneralSecurityException {
        KeyFactory kf = KeyFactory.getInstance("RSA");
        this.privateKey = kf.generatePrivate(
                new PKCS8EncodedKeySpec(decodePem(privPem)));      (*@\co{1}@*)
        this.publicKey = kf.generatePublic(
                new X509EncodedKeySpec(decodePem(pubPem)));
    }

    public String generateAccessToken(User user) {
        return Jwts.builder()
                .header().keyId("2026-09-key1").and()             (*@\co{2}@*)
                .subject(String.valueOf(user.getId()))
                .claim("email", user.getEmail())
                .claim("role", user.getRole().name())
                .issuedAt(new Date())
                .expiration(new Date(System.currentTimeMillis()
                        + accessTokenExpiryMs))
                .signWith(privateKey, Jwts.SIG.RS256)             (*@\co{3}@*)
                .compact();
    }

    private static byte[] decodePem(String pem) {
        String body = pem.replaceAll("-----BEGIN [A-Z ]+-----", "")
                         .replaceAll("-----END [A-Z ]+-----", "")
                         .replaceAll("\\s", "");
        return Base64.getDecoder().decode(body);
    }
}
\end{lstlisting}

\begin{description}
  \item[\co{1}] PKCS\#8 là định dạng chứa mà \code{openssl genpkey} ghi
  cho khóa riêng; X.509 \code{SubjectPublicKeyInfo} là thứ
  \code{-pubout} ghi. Bóc vỏ PEM, giải mã base64, xong.
  \item[\co{2}] \code{kid} (key id) đặt tên cho khóa đã ký token này.
  Tốn vài byte bây giờ và cho phép xoay khóa về sau.
  \item[\co{3}] JJWT suy ra RS256 từ loại khóa; nêu tên tường minh để
  ghi lại ý định.
\end{description}

Gateway --- và bất kỳ service nào muốn tự xác minh --- chỉ giữ khóa
công khai:

\begin{lstlisting}[language=Java]
public JwtAuthenticationFilter(@Value("${app.jwt.public-key-pem}") String pubPem)
        throws GeneralSecurityException {
    this.publicKey = KeyFactory.getInstance("RSA")
            .generatePublic(new X509EncodedKeySpec(decodePem(pubPem)));
}

Claims claims = Jwts.parser()
        .verifyWith(publicKey)      // PublicKey: accepts RS/PS/ES only
        .build()
        .parseSignedClaims(token)
        .getPayload();
\end{lstlisting}

Không gì khác thay đổi: cùng claim, cùng header, cùng
\code{GatewayHeaderAuthFilter}. Gateway không còn đúc được token, và
\code{api-gateway.yml} không cần \code{app.jwt.secret} nữa.

\subsection{Xoay khóa, và \code{kid} để làm gì}

Khóa có tuổi: chính sách bảo xoay mỗi 90 ngày, hoặc một chiếc laptop
chứa file PEM bị mất. Xoay khóa mà không đăng xuất toàn bộ người dùng là
một nghi thức bốn bước:

\begin{enumerate}
  \item Sinh cặp khóa mới; thêm khóa công khai của nó vào JWKS bên cạnh
  khóa cũ, dưới một \code{kid} mới.
  \item Đợi tới khi mọi bên xác minh đã làm mới cache (hoặc khởi động
  lại chúng).
  \item Chuyển auth-service sang ký bằng khóa riêng mới. Token mới mang
  \code{kid} mới; token cũ, còn đang lưu hành, mang \code{kid} cũ --- và
  bên xác minh chọn đúng khóa công khai theo \code{kid}.
  \item Sau khi thời gian sống dài nhất của token đã trôi qua (ở đây 15
  phút cho access token), gỡ khóa cũ khỏi JWKS.
\end{enumerate}

Không có \code{kid}, bên xác minh sẽ phải thử mọi khóa nó biết; có nó,
tra cứu là chính xác và khóa cũ có thể nghỉ hưu một cách chắc chắn. Phiên
bản PEM tĩnh ở trên chưa có endpoint JWKS; bước tự nhiên tiếp theo là
phục vụ \code{/.well-known/jwks.json} từ auth-service và để các bên xác
minh dùng \code{NimbusJwtDecoder.withJwkSetUri(...)} của Spring
Security, thứ tự động cache và tải lại khi gặp \code{kid} lạ.

\section{Token béo phì (payload bloat)}
\label{sec:payload-bloat}

Phi trạng thái rất quyến rũ: nếu token mang mọi thứ một service cần,
không service nào phải gọi service khác để hỏi ``user này được làm việc
này không?'' Các đội theo logic đó cho tới khi token mang vai trò,
quyền, nhóm, feature flag và hồ sơ của user --- và nặng 4\,KB. Phép tính
lúc đó quay lại chống họ. Một trang dashboard bắn 20 lời gọi API song
song gửi đi 80\,KB token trước khi có một byte dữ liệu nghiệp vụ nào.
Nhân với vài chục nghìn user đồng thời và header xác thực chiếm một phần
đo được của băng thông vào. Cookie, nếu token nằm ở đó, có trần 4\,KB;
một số proxy giới hạn kích thước header ở 8\,KB; log ghi header phình
theo.

Access token của dự án mỏng. \code{JwtService.generateAccessToken} ghi
\code{sub}, \code{email}, \code{role}, \code{iat}, \code{exp} --- khoảng
200 byte đã ký. Giữ nó như vậy bằng một quy tắc: \textbf{token mang danh
tính; quyền hạn thì tra cứu.} Khoảnh khắc course-service cần biết ``giáo
viên này được sửa khóa học nào?'', cám dỗ thêm claim \code{courseIds}
xuất hiện. Hãy cưỡng lại. Danh sách đó thay đổi khi một khóa học được
tạo, điều token không thể biết cho tới khi hết hạn; và nó lớn không giới
hạn. Service vốn đã sở hữu dữ liệu ghi danh (Chương~7) --- một truy vấn
có index trả lời câu hỏi đúng và cập nhật.

\section{Mẫu phantom token}

Giả sử token béo là không tránh khỏi --- một mô hình phân quyền cũ, một
tích hợp với đối tác. Mẫu phantom token giữ token \emph{bên ngoài} nhỏ
xíu trong khi token \emph{bên trong} vẫn tự chứa:

\begin{enumerate}
  \item Auth-service tạo JWT đầy đủ nhưng \emph{không} trả về. Nó cất
  JWT vào một kho nhanh (Redis) dưới một định danh mờ ngẫu nhiên, và chỉ
  trả về định danh đó --- một UUID, 36 byte.
  \item Gateway nhận UUID, tra cứu, và thay bằng JWT thật trước khi
  chuyển tiếp. Service nội bộ thấy một token ký bình thường.
  \item Service nội bộ tự xác minh nó (RS256, JWKS). Bên trong vành đai,
  hệ thống phi trạng thái và zero trust; bên ngoài, client mang một token
  không có nội dung nào đọc được.
\end{enumerate}

Đăng xuất trở nên tầm thường và \emph{tức thì}: \code{DEL <uuid>} trong
Redis, và request kế tiếp với định danh đó thất bại tại gateway với 401,
bất kể JWT bên trong còn sống bao lâu. Đây là phần thưởng thật của mẫu
này --- nó sửa điểm yếu cấu trúc duy nhất của JWT, không thu hồi được.

Dự án đã chạy một nửa của nó. Nhìn
\code{JwtService.generateRefreshTokenValue}:

\begin{lstlisting}[language=Java]
public String generateRefreshTokenValue() {
    return UUID.randomUUID().toString();
}
\end{lstlisting}

Refresh token không phải JWT. Nó là một định danh mờ lưu trong bảng
\code{refresh\_tokens} (Flyway \code{V2}) và bị thu hồi bằng cách xóa
dòng --- \code{AuthService.logout} làm đúng thế. Access token vẫn là
một JWT tự chứa. Vậy độ trễ thu hồi của dự án bị chặn trên bởi 15 phút
của access token, và \emph{cơ chế} thu hồi của nó là mẫu phantom áp dụng
cho riêng thông tin xác thực sống dài.

\begin{pitfall}[refresh không xoay vòng]
Đọc \code{AuthService.refreshToken}: nó trả về \emph{đúng} refresh
token nó nhận được. Một refresh token bị đánh cắp một lần vì thế hợp lệ
suốt bảy ngày, và vụ trộm không thể phát hiện. Thực hành trong ngành là
\emph{xoay vòng} (rotation): mỗi lần \code{/refresh} phát một refresh
token mới và vô hiệu token cũ. Nếu token cũ có lúc nào bị đưa ra lại,
hai bên đang giữ bản sao --- thu hồi cả họ token. Bảng đã có đủ cột để
làm việc này; thay đổi là một cặp xóa-và-chèn trong một transaction.
\end{pitfall}

\begin{decision}[ba cách thu hồi một access token]
Khi admin vô hiệu một tài khoản, câu trả lời của dự án là ``trong vòng
15 phút''. Các phương án, theo chi phí tăng dần:
\begin{enumerate}
  \item \textbf{Đợi cho hết hạn.} Access token ngắn cộng refresh token
  có trạng thái. Không thêm hạ tầng; độ trễ thu hồi bằng thời gian sống
  của access token. Chấp nhận được với phần lớn ứng dụng, và là điều dự
  án đang làm.
  \item \textbf{Danh sách cấm (deny-list).} Thêm claim \code{jti} (id
  token); khi đăng xuất hoặc vô hiệu, ghi \code{jti} vào Redis với TTL
  bằng phần đời còn lại của token; gateway tra Redis mỗi request. Một
  lượt tra nhanh mỗi request, và chỉ token bị thu hồi mới được lưu.
  \item \textbf{Phantom token.} Định danh mờ bên ngoài, JWT bên trong,
  Redis làm bảng ánh xạ. Thu hồi tức thì và trình duyệt không có token
  đọc được --- nhưng Redis giờ nằm trên đường găng của \emph{mọi}
  request. Redis chết nghĩa là mọi user bị đăng xuất cùng lúc.
\end{enumerate}
Phản biện của một kỹ sư senior với sự hào hứng của bài viết: phương
án~3 đổi ``không thu hồi được'' lấy ``một điểm lỗi đơn mới''. Phần lớn
đội nên với tới phương án~2 trước và coi phương án~3 là phù hợp với yêu
cầu cấp ngân hàng, nơi chính bài viết cũng thừa nhận session có thể là
công cụ tốt hơn.
\end{decision}

\section{Chuyện gì gãy lúc 3 giờ sáng}

Phân tích sự cố là cách nhanh nhất để phân biệt hai mô hình khóa.
Auth-service chết lúc 3 giờ sáng; không gì khác chết.

\begin{itemize}
  \item \textbf{HS256, như đang triển khai.} Gateway giữ secret và xác
  minh cục bộ: mọi user đang đăng nhập tiếp tục dùng được cho tới khi
  access token hết hạn. Lời gọi refresh thất bại (auth-service là
  endpoint \code{/refresh}), nên phiên chết trong vòng 15 phút. Đăng
  nhập mới thất bại. Hồi phục tức thì khi auth-service trở lại.
  \item \textbf{RS256 với JWKS đã cache.} Giống hệt với các bên xác
  minh đang chạy: chúng đã cache khóa công khai lúc khởi động và không
  cần gì từ auth-service. Sự cố mới là một bên xác minh \emph{lạnh} ---
  một pod gateway khởi động lại trong lúc sự cố không tải được JWKS, và
  tùy cấu hình sẽ từ chối mọi token (fail closed) hoặc crash lúc boot.
  Giảm nhẹ: kèm khóa công khai hiện hành trong cấu hình làm phương án
  dự phòng, hoặc cho cache JWKS một thời gian ân hạn dài.
  \item \textbf{Phantom token.} Thêm Redis vào danh sách những thứ phải
  sống. Auth-service chết có cùng hệ quả như trên; Redis chết là sập
  toàn bộ lưu lượng đã xác thực.
\end{itemize}

Bài học khái quát: phi trạng thái mua cho bạn sự độc lập khỏi bên phát
hành tại thời điểm request. Mọi cơ chế bạn thêm để giành lại quyền kiểm
soát --- JWKS, deny-list, ánh xạ mờ --- đều đưa trở lại một phụ thuộc,
và bạn phải gọi tên được nó và nói được chuyện gì xảy ra khi nó không
sẵn sàng.

\section{Khi JWT là công cụ sai}

Bài viết khép lại bằng một cảnh báo đáng nhắc lại, vì chính con đường
của dự án (Chương~27) có một phương án monolith. JWT xứng đáng với độ
phức tạp của nó khi \emph{nhiều bên độc lập} phải trả lời ``ai đang
gọi?'' mà không chia sẻ kho session: microservices, đăng nhập một lần
xuyên ứng dụng, liên kết OAuth2/OIDC. Nó là công cụ sai khi hệ thống là
một khối triển khai duy nhất có đủ chỗ cho session trong bộ nhớ, khi
quyền thay đổi thường xuyên và phải có hiệu lực ngay, hoặc khi cơ quan
quản lý đòi một phiên phải chấm dứt được trong mili giây. Những trường
hợp đó, session phía server với cookie đơn giản hơn, thu hồi được, và
nhàm chán --- mà nhàm chán là một tính năng.

\section{Ngoài dự án}

Keycloak, Auth0, Cognito, Okta và Entra ID đều phát hành token RS256
hoặc ES256 và công bố JWKS; mọi resource server Spring trong ngành xác
minh bằng cấu hình một dòng \code{jwk-set-uri} đã thấy ở Chương~5.
Service mesh (Istio, Linkerd) hiện thực biến thể mTLS của zero trust để
code ứng dụng vẫn đọc header trong khi tầng vận chuyển chứng minh ai
viết chúng. Mẫu phantom token được sản phẩm hóa bởi Curity và các API
gateway như Kong, Apigee dưới các tên ``token introspection'' và
``opaque token exchange''. Biết phiên bản tự tay làm --- một secret
HS256, hai nơi giữ --- là thứ khiến những sản phẩm đó đọc hiểu được.

\section{Câu hỏi từ phỏng vấn thật}

Những câu này được hỏi như bài toán mở, không phải định nghĩa. Câu trả
lời dưới đây là điều một kỹ sư senior nói trong hai phút đầu; người
phỏng vấn sau đó đào vào đúng đánh đổi bạn vừa nêu.

\subsection*{Q: Làm sao vô hiệu một JWT trước khi nó hết hạn?}

Bạn không làm được; bạn giới hạn thiệt hại. Access token sống vài phút,
không phải vài giờ, nên ``đợi hết hạn'' là mức nền. Muốn thu hồi tức
thì, thêm claim \code{jti} và một deny-list trong Redis với TTL bằng
phần đời còn lại của token --- một lượt tra mỗi request, chỉ lưu token
bị thu hồi, và khi Redis không sẵn sàng có thể fail open hay fail closed
theo chính sách. Mờ hoàn toàn (phantom token) cho thu hồi tức thì nhưng
đặt Redis lên đường găng của mọi request. Trong dự án này refresh token
đã là mờ và lưu trong cơ sở dữ liệu, nên ``đăng xuất mọi nơi'' là
\code{DELETE FROM refresh\_tokens WHERE user\_id = ?} cộng thêm nhiều
nhất 15 phút truy cập còn sót.

\subsection*{Q: Lưu lượng đăng nhập tăng từ 10 lên 10\,000 mỗi giây. Cái gì gãy?}

Ký trước tiên: ký RS256 tốn khoảng 1\,ms CPU mỗi token, nên mười nghìn
lần đăng nhập mỗi giây là mười nhân CPU chỉ để ký, trong khi HS256 gần
như không đáng kể. Rồi tới băm mật khẩu, thứ cố ý chậm (bcrypt cost 10
khoảng 100\,ms) --- đó mới là bức tường thật, và câu trả lời là mở rộng
ngang auth-service cộng rate limit theo IP và theo tài khoản tại
gateway. Xác minh không phải vấn đề: verify RS256 mất vài chục micro
giây và không cần mạng. Cũng kiểm tra rằng lệnh insert refresh token
không dồn lên một index nóng và bảng \code{refresh\_tokens} được dọn các
dòng hết hạn.

\subsection*{Q: Hai region, một hệ thống auth. Cái gì phải chia sẻ?}

Với HS256 là secret, một bài toán phân phối bạn không muốn có xuyên
region. Với RS256 chỉ là bộ khóa công khai, thứ cache được, công khai,
và có thể trễ an toàn trong khoảng chồng lấn của chu kỳ xoay khóa. Đồng
hồ quan trọng hơn khóa: \code{exp} và \code{iat} được so với đồng hồ
của bên xác minh, nên cho phép một độ lệch nhỏ (\code{clockSkewSeconds}
của JJWT, thường 30--60\,s) và chạy NTP ở mọi nơi. Refresh token có
trạng thái, nên hoặc ghim refresh về region phát hành hoặc sao chép
bảng; Postgres đơn lẻ của dự án sẽ là thứ đầu tiên phải đổi.

\subsection*{Q: Một token bị đánh cắp từ trình duyệt của người dùng. Bạn làm gì?}

Khoanh vùng, rồi phát hiện. Khoanh vùng: thu hồi refresh token của
user; access token chết trong vòng đời của nó. Phát hiện: xoay vòng
refresh token biến một token cũ bị dùng lại thành tín hiệu --- hai bên
cùng đưa ra một họ token nghĩa là bị trộm, thu hồi cả họ. Phòng ngừa:
giữ refresh token trong cookie \code{HttpOnly} để script bị tiêm không
đọc được, và giữ access token trong bộ nhớ thay vì
\code{localStorage} (Chương~38). Chỉ thêm ràng buộc thiết bị hoặc IP
nếu bạn chấp nhận chi phí báo động giả với người dùng di động.

\subsection*{Q: Sao không kiểm tra cơ sở dữ liệu ở mọi request cho xong?}

Có thể, và monolith làm thế --- đó là session. Lý do không làm trong
kiến trúc service là mọi service sẽ cần kho user hoặc một lời gọi tới
service có nó, đặt auth-service lên đường găng của mọi thứ và ghép mọi
lần deploy vào nó. Token có chữ ký cho mỗi chặng tự trả lời ``đây là
ai'' tại chỗ. Câu trả lời thật thà nêu cái giá --- độ trễ thu hồi ---
và cái bạn đã trả để mua lại quyền kiểm soát: vòng đời ngắn, refresh
token có trạng thái, và tùy chọn một deny-list.

\subsection*{Q: Trong token của bạn có gì, và cố ý không có gì?}

Subject, vai trò, thời điểm phát hành, hạn dùng, và ở đây một e-mail
--- danh tính, không phải quyền hạn. Không có: danh sách permission,
thành viên nhóm, bất cứ thứ gì thay đổi nhanh hơn hạn token, bất cứ thứ
gì service khác sở hữu, và bất cứ thứ gì bí mật (payload đọc được). Nếu
người phỏng vấn ép ``cho permission vào để khỏi query'', câu trả lời
senior là phép tính ở Mục~\ref{sec:payload-bloat}: 4\,KB nhân mọi
request, và một quyết định phân quyền lỗi thời bị đóng băng suốt đời
token.

\begin{quiz}
\begin{enumerate}
  \item Sau khi chuyển sang RS256, \code{api-gateway.yml} không còn
  chứa \code{app.jwt.secret}. Nếu toàn bộ config-server bị lộ, kẻ tấn
  công có giả mạo được token không? Giải thích file nào mới là thứ phải
  lộ.
  \item Mô tả nghi thức xoay khóa 90 ngày và nêu chính xác \code{kid}
  giải quyết vấn đề gì ở bước~3.
  \item Ba mươi service, mỗi cái xác minh token RS256 bằng JWKS tải lúc
  khởi động. Auth-service chết lúc 3 giờ sáng. User nào bị ảnh hưởng,
  khi nào, và sự kiện duy nhất nào trong lúc sự cố khiến một bên xác
  minh từ chối \emph{toàn bộ} lưu lượng?
  \item Một đồng nghiệp đề xuất thêm claim \code{courseIds} để
  course-service phân quyền mà không cần truy vấn. Nêu hai lý do từ
  chối không liên quan tới kích thước token.
  \item Gateway chuyển tiếp request tới đường dẫn công khai mà không
  đụng header. Giải thích vì sao hôm nay đây là lỗ hổng tiềm ẩn chứ chưa
  phải đang hoạt động, và nêu cấu hình một dòng loại bỏ sự tiềm ẩn đó.
  \item Admin vô hiệu một user đang đăng nhập. Cho biết độ trễ thu hồi
  với thiết kế hiện tại của dự án, và nêu phụ thuộc thêm mà mỗi thiết
  kế nhanh hơn trong hai thiết kế còn lại đưa vào.
\end{enumerate}
\end{quiz}
